%! Two Categorical Variables — Unit 2, Lesson 2.1 — Slides (Beamer)
\documentclass[aspectratio=169, 11pt]{beamer}
\usepackage{apstats-beamer}   % provides \navyheader{} and \sectionlabel[color]{}

% The C and Y column types live in apstats-article, which a beamer deck does not
% load — redeclare them here rather than pulling in the article preamble.
\newcolumntype{C}[1]{>{\centering\arraybackslash}p{#1}}
\newcolumntype{Y}{>{\centering\arraybackslash}X}

\begin{document}

% TITLE SLIDE (CourseName is not defined in beamer, so the name is literal)
{
\setbeamercolor{background canvas}{bg=navy}
\begin{frame}[plain]
  \vspace{2.8cm}\hspace{0.55cm}%
  \begin{minipage}{0.9\textwidth}
    {\color{goldacc}\bfseries\small LESSON 2.1}\\[0.5cm]
    {\color{white}\bfseries\LARGE Two Categorical Variables}\\[1.2cm]
    {\color{skymid}\small Statistics~$\cdot$~Unit 2}
  \end{minipage}
\end{frame}
}

% GRADUAL RELEASE deck: targets → warm-up → I DO divider → one frame per notes row (Two-way
% table, Conditional, Association) → WE DO (Texting drivers) → spoken debrief → close & START
% THE HOMEWORK, which is the individual practice. No group activity, no exit ticket.
\newcommand{\redink}[1]{{\color{redacc}\bfseries #1}}

% ── LEARNING TARGETS ─────────────────────────────────────────────────────────
\begin{frame}[t, plain]
  \navyheader{Today's targets}
  \sectionlabel{Lesson 2.1}
  By the end of today, I can\ldots
  \begin{itemize}\setlength{\itemsep}{5pt}
    \item ask whether two variables are related, and read a \textbf{two-way table}.
    \item find \textbf{joint}, \textbf{marginal}, and \textbf{conditional relative
          frequencies} --- dividing by the right total.
    \item read a \textbf{segmented bar graph} and decide whether there is an
          \textbf{association}.
    \item explain why an association in a survey is not proof of cause.
  \end{itemize}
  \begin{block}{How today works}
    Warm-up $\to$ \textbf{notes together} $\to$ \textbf{one we work as a class} $\to$
    debrief $\to$ \textbf{you start the homework here, alone}.
  \end{block}
\end{frame}

% ── WARM-UP ──────────────────────────────────────────────────────────────────
\begin{frame}[t, plain]
  \navyheader{Warm-Up}
  \sectionlabel{5 minutes $\cdot$ on your own}
  \begin{enumerate}\setlength{\itemsep}{8pt}
    \item 45 of 180 students walk to school. Relative frequency?
    \item Section A: 12 of 20 passed. Section B: 18 of 40 passed. \\
          Which section had \emph{more} students pass? Which section did \emph{better}?
    \item Phone in bedroom (yes/no) $\cdot$ hours of sleep $\cdot$ plays a sport (yes/no): \\
          categorical or quantitative? Write a question asking whether two are
          \emph{related}.
  \end{enumerate}
  \vspace{6pt}
  \begin{center}
    {\small\color{slate} Keep item 2 in mind: \redink{more} and \redink{better} were different
    answers.}
  \end{center}
\end{frame}

% ── I DO — direct instruction begins ─────────────────────────────────────────
{
\setbeamercolor{background canvas}{bg=navy}
\begin{frame}[plain]
  \vspace{1.8cm}\hspace{0.8cm}%
  \begin{minipage}{0.86\textwidth}
    {\color{goldacc}\bfseries\small GUIDED NOTES \ \textbullet\ \ I DO}\\[0.7cm]
    {\color{white}\bfseries\Large Open to your Guided Notes page. Fill each blank as we
    name it --- not before.}\\[0.9cm]
    {\color{skymid}\small Three terms go in the vocabulary box today: two-way table,
    conditional relative frequency, association.}
  \end{minipage}
\end{frame}
}

% ── ROW 1: TWO-WAY TABLE ─────────────────────────────────────────────────────
\begin{frame}[t, plain]
  \navyheader{Reading a two-way table}
  \sectionlabel{notes, Two-Way Table $\cdot$ problems 1--3}
  200 students: does your phone stay in your bedroom overnight? Do you sleep 8+ hours?
  \vspace{6pt}
  \begin{center}
  \renewcommand{\arraystretch}{1.25}
  \begin{tabular}{l | c c | c}
     & \textbf{8+ hours} & \textbf{Under 8} & \textbf{Total} \\ \hline
    \textbf{Phone in room} & 48 & 112 & 160 \\
    \textbf{No phone} & 20 & 20 & 40 \\ \hline
    \textbf{Total} & 68 & 132 & \textbf{200} \\
  \end{tabular}
  \end{center}
  \vspace{6pt}
  \begin{columns}[t]
    \begin{column}{0.46\textwidth}
      \begin{block}{Joint relative frequency}
        One cell $\div$ the \textbf{grand total}: \\
        phone \emph{and} 8+ $= 48/200 = 0.24$ (UNC-1.P.4)
      \end{block}
    \end{column}
    \begin{column}{0.46\textwidth}
      \begin{block}{Marginal relative frequency}
        A row or column total $\div$ the \textbf{grand total}: \\
        8+ hours $= 68/200 = 0.34$ (UNC-1.Q.1)
      \end{block}
    \end{column}
  \end{columns}
\end{frame}

% ── ROW 2: CONDITIONAL ───────────────────────────────────────────────────────
\begin{frame}[t, plain]
  \navyheader{``Of those who\ldots'' picks the denominator}
  \sectionlabel{notes, Conditional $\cdot$ problems 4--6}
  \vspace{4pt}
  \begin{block}{Conditional relative frequency (UNC-1.Q.2)}
    Look inside \textbf{one} row or column only, and divide by \textbf{that group's total}.
  \end{block}
  \vspace{8pt}
  \begin{columns}[t]
    \begin{column}{0.46\textwidth}
      \textbf{Of phone-in-room students}, 8+ hours: \\[3pt]
      \hspace{1em}$48 / \mathbf{160} = 0.30$ \\[8pt]
      \textbf{Of no-phone students}, 8+ hours: \\[3pt]
      \hspace{1em}$20 / \mathbf{40} = 0.50$
    \end{column}
    \begin{column}{0.46\textwidth}
      \textbf{Of 8+ hour sleepers}, phone in room: \\[3pt]
      \hspace{1em}$48 / \mathbf{68} \approx 0.71$ \\[8pt]
      {\small\color{slate} Same 48 on top. A \redink{different group} on the bottom --- so a
      different question and a different answer.}
    \end{column}
  \end{columns}
\end{frame}

% ── ROW 3: ASSOCIATION — THE CRUX ────────────────────────────────────────────
\begin{frame}[t, plain]
  \navyheader{48 is bigger than 20. So what?}
  \sectionlabel{notes, Association $\cdot$ problems 7--10 $\cdot$ the whole point of today}
  \begin{columns}[c]
    \begin{column}{0.36\textwidth}
      \begin{center}
      \begin{tikzpicture}[x=1.0cm, y=0.030cm]
        \fill[navy]    (0.3,0)  rectangle (1.3,30);
        \fill[goldacc] (0.3,30) rectangle (1.3,100);
        \fill[navy]    (1.8,0)  rectangle (2.8,50);
        \fill[goldacc] (1.8,50) rectangle (2.8,100);
        \draw (0.3,0) rectangle (1.3,100); \draw (1.8,0) rectangle (2.8,100);
        \node[white, font=\scriptsize\bfseries] at (0.8,15) {30\%};
        \node[font=\scriptsize\bfseries] at (0.8,65) {70\%};
        \node[white, font=\scriptsize\bfseries] at (2.3,25) {50\%};
        \node[font=\scriptsize\bfseries] at (2.3,75) {50\%};
        \draw[thick] (0,0) -- (3.0,0); \draw[thick] (0,0) -- (0,104);
        \node[below, font=\tiny, align=center] at (0.8,-2) {Phone in\\room (160)};
        \node[below, font=\tiny, align=center] at (2.3,-2) {No phone\\(40)};
        \fill[navy] (0.3,110) rectangle (0.5,118); \node[right, font=\tiny] at (0.5,114) {8+ hours};
        \fill[goldacc] (1.8,110) rectangle (2.0,118); \node[right, font=\tiny] at (2.0,114) {Under 8};
      \end{tikzpicture}
      \end{center}
    \end{column}
    \begin{column}{0.60\textwidth}
      \small
      \begin{itemize}\setlength{\itemsep}{4pt}
        \item ``48 phone students sleep 8+, only 20 no-phone students do.'' True --- and
              \redink{misleading}: there are 160 phone students and only 40 without.
        \item Compare \textbf{percents}: 30\% vs.\ 50\%. The bars differ, so the variables are
              \textbf{associated} (UNC-1.R.1).
        \item A \textbf{mosaic plot} would also make each bar's \emph{width} show its group's
              size.
      \end{itemize}
    \end{column}
  \end{columns}
  \vspace{2pt}
  \begin{block}{Association is not causation}
    This is a survey. Students who keep phones out may differ in other ways --- the pattern
    could even be chance (VAR-1.D.1). It does not show that phones \emph{cause} less sleep.
  \end{block}
\end{frame}

% ── WE DO ────────────────────────────────────────────────────────────────────
\begin{frame}[t, plain]
  \navyheader{Guided Practice: Texting drivers}
  \sectionlabel[goldacc]{We do $\cdot$ problems 11--14 $\cdot$ you hold the pen}
  \begin{columns}[t]
    \begin{column}{0.46\textwidth}
      \begin{center}
      \renewcommand{\arraystretch}{1.25}
      \begin{tabular}{l | c c | c}
         & \textbf{Texted} & \textbf{Did not} & \textbf{Total} \\ \hline
        \textbf{16--24} & 20 & 30 & 50 \\
        \textbf{25+} & 36 & 114 & 150 \\ \hline
        \textbf{Total} & 56 & 144 & 200 \\
      \end{tabular}
      \end{center}
    \end{column}
    \begin{column}{0.50\textwidth}
      \begin{itemize}\setlength{\itemsep}{5pt}
        \item What proportion of all drivers texted? Joint or marginal?
        \item What proportion texted in \emph{each} age group?
        \item ``Older drivers text more --- 36 vs.\ 20.'' Respond.
        \item Associated? Say it in context.
      \end{itemize}
    \end{column}
  \end{columns}
  \vspace{8pt}
  \begin{block}{The question behind every answer}
    \emph{Of whom?} Name the group before you divide --- and compare \redink{percents}, never
    counts, when the groups are different sizes.
  \end{block}
\end{frame}

% ── DEBRIEF ──────────────────────────────────────────────────────────────────
\begin{frame}[t, plain]
  \navyheader{Debrief: more is not the same as more likely}
  \sectionlabel{8 minutes $\cdot$ pens down, papers up --- we say these aloud}
  \vspace{6pt}
  \renewcommand{\arraystretch}{1.3}
  \begin{tabularx}{\textwidth}{l Y Y}
    & \textbf{Count} & \textbf{Percent of the group} \\ \hline
    Phone in room, 8+ hours & 48 & 30\% \\
    No phone, 8+ hours & 20 & \textbf{50\%} \\
    Drivers 25+, texted & 36 & 24\% \\
    Drivers 16--24, texted & 20 & \textbf{40\%} \\
  \end{tabularx}
  \vspace{10pt}
  \begin{block}{Say it out loud}
    The bigger count came from the \redink{bigger group}. Association lives in the conditional
    percents --- and even then, a survey cannot tell us \emph{why}.
  \end{block}
\end{frame}

% ── CLOSE & START THE HOMEWORK (the individual practice, scored) ────────────
\begin{frame}[t, plain]
  \navyheader{Homework --- start it now}
  \sectionlabel[goldacc]{5 items $\cdot$ scored $\cdot$ due the first class after two study halls}
  A gym's 250 newest members: monthly or annual plan, still attending or stopped.
  \begin{itemize}\setlength{\itemsep}{4pt}
    \item \textbf{Part A:} a joint relative frequency, and one ``of those who \emph{stopped}.''
    \item \textbf{Part B:} each plan's attendance rate, the manager's claim, and one
          multiple-choice item with a justification.
  \end{itemize}
  \begin{block}{Silent and alone --- I am circulating}
    \emph{A2 and B2} are the ones I am reading over your shoulder. The AP Practice page before
    it is \textbf{not scored} --- work it for the reps and bring questions to study hall.
  \end{block}
\end{frame}

% ── CLOSE ────────────────────────────────────────────────────────────────────
{
\setbeamercolor{background canvas}{bg=navy}
\begin{frame}[plain]
  \vspace{1.5cm}\hspace{0.8cm}%
  \begin{minipage}{0.86\textwidth}
    {\color{goldacc}\bfseries\small BEFORE YOU GO}\\[0.7cm]
    {\color{white}\bfseries\Large Ask ``of whom?'' before you divide. Compare percents, not
    counts.}\\[0.9cm]
    {\color{skymid}\small Next: Lesson 2.2 --- Scatterplots and Correlation, when both
    variables are numbers.}
  \end{minipage}
\end{frame}
}

\end{document}
